% ========================================
% 通用教材LaTeX基础配置模板
% 版本: v2.0
% 适用: 大学教材编写
% 更新: 2025年8月7日
% ========================================

% 文档类配置 - 中文书籍类
\documentclass[
    12pt,           % 字体大小
    a4paper,        % 纸张大小
    twoside,        % 双面排版
    openright       % 章节从右页开始
]{ctexbook}

% ========================================
% 基础包引入
% ========================================

% 数学支持
\usepackage{amsmath,amssymb,amsthm}

% 图形和颜色
\usepackage{graphicx}
\usepackage{xcolor}
\usepackage{tikz}
\usetikzlibrary{shapes,arrows,positioning}

% 页面布局
\usepackage[
    top=2.5cm,
    bottom=2.5cm,
    left=2.8cm,
    right=2.2cm,
    bindingoffset=0.5cm,
    headheight=15pt,
    footskip=1.5cm
]{geometry}

% 字体和编码
\usepackage{fontspec}
\usepackage{setspace}
\onehalfspacing  % 1.5倍行距

% 表格增强
\usepackage{longtable}
\usepackage{booktabs}
\usepackage{array}
\usepackage{multirow}
\usepackage{multicol}

% 列表环境
\usepackage{enumitem}

% 代码环境
\usepackage{listings}

% 美化框架
\usepackage{tcolorbox}
\tcbuselibrary{most}

% 页眉页脚
\usepackage{fancyhdr}

% 超链接
\usepackage{hyperref}

% 参考文献
\usepackage{cite}

% 图标字体
\usepackage{fontawesome5}

% ========================================
% 中文字体配置
% ========================================

% 中文字体配置 - Windows系统优化
\setCJKmainfont{SimSun}[
    BoldFont=SimHei,
    ItalicFont=KaiTi
]
\setCJKsansfont{SimHei}
\setCJKmonofont{FangSong}

% 英文字体配置
\setmainfont{Times New Roman}
\setsansfont{Arial}
\setmonofont{Courier New}

% ========================================
% 配色方案定义
% ========================================

% 主色调 - 可根据学科调整
\definecolor{primarycolor}{RGB}{52, 152, 219}      % 主蓝色
\definecolor{secondarycolor}{RGB}{230, 126, 34}    % 橙色
\definecolor{accentcolor}{RGB}{39, 174, 96}        % 绿色
\definecolor{warningcolor}{RGB}{241, 196, 15}      % 黄色
\definecolor{dangercolor}{RGB}{231, 76, 60}        % 红色

% 灰度系列
\definecolor{lightgray}{RGB}{248, 249, 250}
\definecolor{mediumgray}{RGB}{173, 181, 189}
\definecolor{darkgray}{RGB}{73, 80, 87}

% 代码配色
\definecolor{codeblue}{RGB}{40, 90, 200}
\definecolor{codegreen}{RGB}{0, 150, 0}
\definecolor{codered}{RGB}{200, 0, 0}
\definecolor{codegray}{RGB}{128, 128, 128}
\definecolor{backcolour}{RGB}{248, 248, 248}

% ========================================
% 页面样式配置
% ========================================

% 页眉页脚样式
\pagestyle{fancy}
\fancyhf{}
\fancyhead[LE]{\leftmark}
\fancyhead[RO]{\rightmark}
\fancyfoot[LE,RO]{\thepage}
\renewcommand{\headrulewidth}{0.4pt}
\renewcommand{\footrulewidth}{0pt}

% 章节样式
\ctexset{
    chapter={
        format={\centering\Huge\bfseries},
        name={第,章},
        number={\arabic{chapter}},  % 使用阿拉伯数字
        beforeskip=20pt,
        afterskip=40pt
    },
    section={
        format={\Large\bfseries},
        beforeskip=20pt,
        afterskip=10pt
    },
    subsection={
        format={\large\bfseries},
        beforeskip=15pt,
        afterskip=8pt
    }
}

% ========================================
% 定理环境配置
% ========================================

% 定义定理样式
\theoremstyle{definition}
\newtheorem{definition}{定义}[chapter]
\newtheorem{theorem}{定理}[chapter]
\newtheorem{lemma}{引理}[chapter]
\newtheorem{example}{例}[chapter]
\newtheorem{exercise}{练习}[chapter]

% ========================================
% 列表样式配置
% ========================================

% 设置列表样式
\setlist[itemize,1]{label=\textbullet}
\setlist[itemize,2]{label=\textendash}
\setlist[itemize,3]{label=\textasteriskcentered}

\setlist[enumerate,1]{label=\arabic*.}
\setlist[enumerate,2]{label=(\alph*)}
\setlist[enumerate,3]{label=\roman*.}

% ========================================
% 标题页配置
% ========================================

% 自定义标题页命令
\newcommand{\makecustomtitle}[6]{
    \begin{titlepage}
        \centering
        \vspace*{2cm}
        
        % 学校logo和名称
        {\Large \textbf{#1} \par}
        \vspace{1cm}
        
        % 教材标题
        {\Huge \textbf{#2} \par}
        \vspace{0.5cm}
        {\Large #3 \par}
        \vspace{2cm}
        
        % 作者信息
        {\large
        \begin{tabular}{rl}
            主编: & #4 \\
            副主编: & #5 \\
        \end{tabular}
        \par}
        
        \vfill
        
        % 出版信息
        {\large #6 \par}
        {\large \today \par}
    \end{titlepage}
}

% ========================================
% 超链接配置
% ========================================

\hypersetup{
    colorlinks=true,
    linkcolor=primarycolor,
    filecolor=primarycolor,
    urlcolor=primarycolor,
    citecolor=primarycolor,
    pdftitle={教材标题},
    pdfauthor={作者姓名},
    pdfsubject={教材主题},
    pdfkeywords={关键词1, 关键词2, 关键词3},
    bookmarksnumbered=true,
    bookmarksopen=true
}

% ========================================
% 自定义命令
% ========================================

% 强调命令
\newcommand{\highlight}[1]{\textcolor{primarycolor}{\textbf{#1}}}
\newcommand{\important}[1]{\textcolor{dangercolor}{\textbf{#1}}}
\newcommand{\note}[1]{\textcolor{secondarycolor}{\textit{#1}}}

% 代码相关命令
\newcommand{\code}[1]{\texttt{#1}}
\newcommand{\file}[1]{\texttt{\textbf{#1}}}
\newcommand{\variable}[1]{\texttt{\textit{#1}}}

% 快捷符号
\providecommand{\checkmark}{\textcolor{accentcolor}{\faCheck}}
\newcommand{\crossmark}{\textcolor{dangercolor}{\faTimes}}
\newcommand{\infomark}{\textcolor{primarycolor}{\faInfoCircle}}
\newcommand{\warningmark}{\textcolor{warningcolor}{\faExclamationTriangle}}

% ========================================
% 结构化文档宏
% ========================================

% 章节摘要
\newcommand{\chaptersummary}[1]{
    \begin{tcolorbox}[
        colback=lightgray,
        colframe=primarycolor,
        title=本章要点,
        fonttitle=\bfseries,
        left=5pt,
        right=5pt,
        top=5pt,
        bottom=5pt
    ]
    #1
    \end{tcolorbox}
}

% 学习目标
\newcommand{\learningobjectives}[1]{
    \begin{tcolorbox}[
        colback=primarycolor!10,
        colframe=primarycolor,
        title=\faLightbulb\ 学习目标,
        fonttitle=\bfseries,
        left=5pt,
        right=5pt,
        top=5pt,
        bottom=5pt
    ]
    #1
    \end{tcolorbox}
}

% 关键概念
\newcommand{\keypoints}[1]{
    \begin{tcolorbox}[
        colback=secondarycolor!10,
        colframe=secondarycolor,
        title=\faKey\ 关键概念,
        fonttitle=\bfseries,
        left=5pt,
        right=5pt,
        top=5pt,
        bottom=5pt
    ]
    #1
    \end{tcolorbox}
}

% 实践练习
\newcommand{\practiceexercise}[1]{
    \begin{tcolorbox}[
        colback=accentcolor!10,
        colframe=accentcolor,
        title=\faCogs\ 实践练习,
        fonttitle=\bfseries,
        left=5pt,
        right=5pt,
        top=5pt,
        bottom=5pt
    ]
    #1
    \end{tcolorbox}
}

% 注意事项
\newcommand{\attention}[1]{
    \begin{tcolorbox}[
        colback=warningcolor!10,
        colframe=warningcolor,
        title=\faExclamationTriangle\ 注意,
        fonttitle=\bfseries,
        left=5pt,
        right=5pt,
        top=5pt,
        bottom=5pt
    ]
    #1
    \end{tcolorbox}
}

% ========================================
% 文档开始前的最后配置
% ========================================

% 设置段落缩进
\setlength{\parindent}{2em}
\setlength{\parskip}{0.5ex plus 0.2ex minus 0.1ex}

% 设置浮动体参数
\renewcommand{\textfraction}{0.15}
\renewcommand{\topfraction}{0.85}
\renewcommand{\bottomfraction}{0.65}
\renewcommand{\floatpagefraction}{0.60}

% 图表标题样式
\usepackage{caption}
\captionsetup{
    font=small,
    labelfont=bf,
    textfont=it,
    margin=10pt
}

% ========================================
% 模板使用说明
% ========================================

% 在文档开始时使用以下命令：
% 1. \makecustomtitle{学校名}{教材标题}{副标题}{主编}{副主编}{出版社}
% 2. \tableofcontents
% 3. 开始章节内容

% 可用的特殊环境：
% - \learningobjectives{内容} - 学习目标框
% - \keypoints{内容} - 关键概念框  
% - \practiceexercise{内容} - 实践练习框
% - \attention{内容} - 注意事项框
% - \chaptersummary{内容} - 章节摘要框

% 常用命令：
% - \highlight{文本} - 高亮强调
% - \important{文本} - 重要内容
% - \note{文本} - 注释说明
% - \code{代码} - 行内代码
% - \checkmark, \crossmark - 检查标记
